\section*{Appendix}

\appendix

\section{Block Cholesky factorization of the augmented Gram matrix}\label{appendix:chol_aug_matrix}
Let $A$ be an $m \times n$ full row rank matrix for which the Cholesky factor of the Gram matrix $AA^\top$ is available. Considering $\calA = \left\{i_1,\ldots, i_p\right\} \subseteq \{1,\ldots,n\}$, let $Z \in \RR^{p\times n}$ be the matrix whose row $k$ is the row $i_k$ of the $n\times n$ identity.
In this appendix, we show how to construct the Cholesky decomposition of the Gram matrix $MM^\top$, where $M$ is defined as
\[M = \begin{bmatrix}
	A \\ Z
\end{bmatrix} \in \RR^{(m+p)\times n}.\]  
Following the reasoning in~\cite[][section 4.2.9]{golubvanloan:2013} for the Cholesky decomposition of a general block matrix, we will show how, at the cost of extra computations, we can recover the Cholesky factor of $MM^\top$ using its block structure and the factor $L$. The definition of $M$ naturally leads to the block pattern
\[MM^\top = \begin{bmatrix}
	AA^\top & AZ^\top \\ ZA^\top & ZZ^\top
\end{bmatrix}\]
Simple computations first show that 
\[ZZ^\top = I_p,\quad AZ^\top = [A_{i_1},\ldots,A_{i_p}].\]
The latter, written $A_{|\calA}$, is the restriction of $A$ to the columns corresponding to the indices in $\calA$. We can write $MM^\top$ as 
\begin{equation}\label{eq:augmented_grammat}
	\begin{bmatrix}
		AA^\top &  A_{|\calA} \\  A_{|\calA}^\top & I_p
	\end{bmatrix}
\end{equation}
Let $\tilde{L}$ be $(m+p)\times(m+p)$ triangular such that $MM^\top=\tilde{L}\tilde{L}^\top$. We apply the same block structure to $\tilde{L}$, i.e.
\[\tilde{L} = \begin{bmatrix}
	\tilde{L}_{11} & 0 \\ G_{21} & \tilde{L}_{22},
\end{bmatrix},\]
with 
\begin{itemize}
	\item $\tilde{L}_{11} \ m\times m$ lower triangular
	\item $G_{21} \ p\times m$
	\item $\tilde{L}_{22} \ p \times p$ lower triangular
\end{itemize}
Verifying $MM^\top= \tilde{L}\tilde{L}^\top$ implies the matrix equality
\begin{equation}\label{eq:block_equalities_chol_fact}
	\begin{bmatrix}
		AA^\top &  A_{|\calA} \\  A_{|\calA}^\top & I_p
	\end{bmatrix}
	=
	\begin{bmatrix}
		\tilde{L}_{11}\tilde{L}_{11}^\top & \tilde{L}_{11}G_{21}^\top \\ 
		G_{21}\tilde{L}_{11}^\top & G_{21}G_{21}^\top + \tilde{L}_{22}\tilde{L}_{22}^\top
	\end{bmatrix}.
\end{equation}
By comparison of the blocks in~\eqref{eq:block_equalities_chol_fact}, it follows that
\begin{equation*}
	\begin{aligned}
		AA^\top &= \tilde{L}_{11}\tilde{L}_{11}^\top \\
		A_{|\calA} &= \tilde{L}_{11}G_{21}^\top \\
		I_p &= G_{21}G_{21}^\top + \tilde{L}_{22}\tilde{L}_{22}^\top.
	\end{aligned}
\end{equation*}
Therefore, we can form $\tilde{L}$ by first setting $\tilde{L}_{11}=L$, then solve $p$ lower triangular systems to compute $G_{21}$ and finally compute the Cholesky factor of $I_p-G_{21}G_{21}^\top$ to get $\tilde{L}_{22}$. We briefly justify why the latter is well defined. By definition of $G_{21}^\top$ and the full rank assumption of $A$, one has
\[G_{21}^\top = \left(\tilde{L}_{11}\right)^{-1}A_{|\calA} = L^{-1}A_{|\calA}.\]
Then, we obtain
\begin{equation*}
	\begin{aligned}
		I_p - G_{21}G_{21}^\top &= I_p - \left(L^{-1}A_{|\calA}\right)^\top\left(L^{-1}A_{|\calA}\right) \\
		&= I_p - A_{|\calA}^\top \left(LL^\top\right)^{-1}A_{|\calA} \\
		&= I_p - A_{|\calA}^\top \left(AA^\top\right)^{-1}A_{|\calA},
	\end{aligned}
\end{equation*}
which is nothing but the Schur complement of $AA^\top$ of the augmented matrix~\eqref{eq:augmented_grammat}. Thus, the positive definiteness of $I_p-G_{21}G_{21}^\top$ comes from the positive definiteness of $MM^\top$ (also implied by the full rank assumption of $A$).

\section{Cauchy point computation}\label{appendix:cauchy_point_computation}

We now describe a procedure to compute the Cauchy point as the first local minimizer of the quadratic model along the projected gradient path adapted from~\cite{conn-etal:1988a} to the case where linear equalities are present. For the sake of clarity, we omit the iteration index during the rest of this subsection. 
The piece-wise arc $s(t) = \proj{\Omega}{x-tg}-x$ satisfies the constraints
\begin{itemize}
	\item $As(t) = 0$
	\item $s^{(l)} \le s(t) \le s^{(u)}$,
\end{itemize}
with $s^{(l)}_i = \max(-\Delta,l_i-x_i)$ and $s^{(u)}_i = \min(\Delta,u_i-x_i)$.

The following breakpoints:
\begin{equation*}
	0=t_0 < t_1 < \ldots < t_p,
\end{equation*} 
correspond to the successive scalars at which at least one component of $s(t)$ satisfies a bound with equality. Note that since $A$ has $m$ rows and is full rank, there are $n-m$ degrees of freedom remaining  and thus at most $n-m$ breakpoints before the projected gradient path is constant. We recall the recursive expression
\[s(t) = -(t-t_i)z_i + s(t_i),\]
for $t\in [t_i,t_{i+1})$ and with the reduced gradient $z_i$ given by~\eqref{eq:reduced_gradient_proj_grad_path}.

To find the first local minimum of the scalar function $\psi(t) = q(s(t))$, we successively study each interval $[t_i,t_{i+1})$ to assert whether or not it contains a local minimizer.
Assume we have not found a local minimizer on $[t_0,t_i)$ and thus look at the interval $[t_i,t_{i+1})$. The model along this arc can be written
\begin{equation}\label{eq:model_projected_gradient_interval}
	\psi(t) = \dfrac{\psi_i''}{2}(\Delta t)^2 + \psi_i'\Delta t
\end{equation}
with
\begin{itemize}
	\item $\psi_i'' = \inner{z_i}{Hz_i}$
	\item $\psi_i' = -\inner{s(t_i)}{Hz_i} - \inner{g}{z_i}$
	\item $\Delta t = t-t_i$
\end{itemize}
Different cases can occur depending on the values of the slope $\psi_i'$ and the curvature $\psi_i''$. 
Firstly, if
\begin{align*}
	&\psi_i' > 0 \text{ or } \\
	&\psi'_i=0 \text{ and } \psi_i'' > 0,
\end{align*}
then $t_i$ is the required minimizer. Next, if
\[\psi_i' < 0 \text{ and } \psi_i'' > 0,\]
the quadratic~\eqref{eq:model_projected_gradient_interval} has a strict minimizer at
\[t_i- \dfrac{\psi_i'}{\psi_i''}.\]
If the latter belongs to the interval of interest, i.e. 
\[t_i-\psi_i'/\psi_i'' < t_{i+1},\]
then it is the required minimizer. In other cases, the minimizer is at or beyond $t_{i+1}$. To prepare for the study of the next interval, we first need to find the next breakpoint, given by the smallest scalar $t_{i+1}>t_i$ such that, at $s(t_{i+1})$, one of the free components hits one of its bounds. By introducing
\[\delta_i^- = \min_{(z_i)_j > 0} \left\{ (s(t_i)_j - s^{(l)}_j)/(z_i)_j\right\}, \qquad \delta_i^+ = \min_{(z_i)_j < 0} \left\{(s(t_i)_j - s^{(u)}_j)/(z_i)_j\right\},\]
the next breakpoint is given by 
\begin{equation}\label{eq:next_breakpoint}
	t_{i+1} = t_i + \delta_i,
\end{equation} 
with $\delta_i = \min\left(\delta_i^-,\delta_i^+\right)$. We then add the corresponding variable index to the list of fixed components, compute the next reduced gradient $z_{i+1}$ and finally update the slope and curvature of the model along the next interval. A last termination case can occur. Since $A$ is of rank $m$, at most $n-m$ bounds can become active so if we never find a local minimum, the procedure still ends whenever it reaches the last breakpoint $t_{p}$. Indeed, past this breakpoint, the model along the projected gradient path is constant so we can return the last accumulated step $s(t_p)$ as the Cauchy step. Note that, in this case, it is also the total step, because no more variables can be modified. The presence of the trust region constraint ensures that all bounds become active. 
The procedure for the Cauchy step computation is outlined in Algorithm~\ref{algo:cauchy_point}.
\begin{algorithm}
	\caption{Cauchy step computation}\label{algo:cauchy_point}
	\begin{algorithmic}[1]
		\Require Hessian $H$, gradient $g$, bounds on the direction $s^{(l)},\ s^{(u)}$
		\State Identify $I(0)$ and compute the direction $z_0$
		\State Set \textbf{found} $\gets$ \textbf{false} 
		\State Set $t_0\gets 0$, $s^{(0)}\gets 0$ and initialize counter $i\gets 0$
		\Repeat
		\State $\psi_i' \gets -\inner{s^{(i)}}{Hz_i}-\inner{g}{z_i},\ \psi_i'' \gets \inner{z_i}{Hz_i}$
		\State Find the next breakpoint $t_{i+1}$ by~\eqref{eq:next_breakpoint}
		\If{$\psi_i' > 0 \text{ or } \psi'_i=0 \text{ and } \psi_i'' > 0$}
		\State Set $s^C \gets s^{(i)}$
		\State \textbf{found $\gets$ \textbf{true}}
		\ElsIf{$\psi_i' < 0 \text{ and } \psi_i'' > 0$ \text{ and } $t_i- \psi_i'/\psi_i'' < t_{i+1}$}
		\State Set $\Delta t \gets - \psi_i'/\psi_i''$
		\State Set$s^C \gets s^{(i)} - \Delta t z_i$, \textbf{found $\gets$ \textbf{true}}
		\Else
		\State Set $\Delta t \gets t_{i+1}-t_i$
		\State Compute the next direction $z_{i+1}$
		\State Set $s^{(i+1)} \gets s^{(i)} - \Delta t z_i$
		\EndIf
		\State Increment $i \gets i+1$
		\If{$|I(t_i)| = n-m$}
		\State Set $s^C \gets s^{(i)}$, \textbf{found} $\gets$ \textbf{true}
		\EndIf
		\Until{\textbf{found}}
		\State \Return $s^C$
	\end{algorithmic}
\end{algorithm}
